\documentclass[11pt, a4paper]{amsart}
\usepackage[utf8]{inputenc}
\usepackage{amsmath, amssymb, amsthm, amsfonts}
\usepackage{mathtools}
\usepackage{bm}
\usepackage[margin=1in]{geometry}
\usepackage{hyperref}

\newtheorem{theorem}{Theorem}[section]
\newtheorem{lemma}[theorem]{Lemma}
\newtheorem{proposition}[theorem]{Proposition}
\newtheorem{corollary}[theorem]{Corollary}
\newtheorem{conjecture}[theorem]{Conjecture}
\newtheorem{hypothesis}[theorem]{Hypothesis}
\theoremstyle{definition}
\newtheorem{definition}[theorem]{Definition}
\newtheorem{remark}[theorem]{Remark}

\newcommand{\R}{\mathbb{R}}
\newcommand{\N}{\mathbb{N}}
\newcommand{\Q}{\mathbb{Q}}
\newcommand{\C}{\mathbb{C}}
\newcommand{\Z}{\mathbb{Z}}
\newcommand{\CP}{\mathbb{CP}}
\newcommand{\HH}{\mathcal{H}}
\newcommand{\pphi}{\varphi}

\title{Cascade $\sigma$-Transport for the Hodge Conjecture:\\
  A Conditional Reduction under Hypothesis $\mathbf{H_{4a}}$}
\author{}
\date{}

\begin{document}

\maketitle

\begin{abstract}
We develop a \emph{conditional $\sigma$-transport reduction} of the
Hodge Conjecture within the cascade-correspondence programme. The
reduction is conditional on Hypothesis $\mathbf{H_{4a}}$ of the
cascade-correspondence foundations
\cite[\S H${}_{4a}$]{CascadeFoundations}, which asserts that every
smooth projective complex variety $X$ admits a cascade-compatible
embedding $\iota_X$ such that the Hodge-class structure on
$H^*(X, \Q)$ is compatible with the $\sigma$-fixed-part projector
$\Pi_\sigma$ on $H^*(X, \Q(\varphi))$ under pullback.

Under $\mathbf{H_{4a}}$, our main theorem
(Theorem~\ref{thm:main-cond}) reduces the Hodge Conjecture for a
general smooth projective complex variety $X$ to the
\emph{algebraicity of the $\sigma$-fixed part}: showing that the
image of $\Pi_\sigma$ on $H^{p,p}(X, \Q(\varphi))$ is spanned by
algebraic cycle classes. We do \emph{not} claim a proof of the
Hodge Conjecture; the algebraicity-of-fixed-part problem is the
remaining classical content.

We discuss five classes of varieties in this paper. \emph{Three
of these classes have classical Hodge proofs available}
(projective spaces $\CP^n$ via Lefschetz $(1,1)$ and cell
decomposition; flag varieties $G/P$ via Schubert decomposition;
smooth complete toric varieties via Danilov's theorem). \emph{The
other two classes are cited only at weaker classical status}: for
CM abelian varieties Deligne proves absolute-Hodge (not Hodge in
general) and Cattani--Deligne--Kaplan prove Hodge-locus
algebraicity (not Hodge for individual classes); for low-codimension
complete intersections, weak Lefschetz handles non-middle degrees,
while Noether--Lefschetz constrains middle-degree primitive classes
only in restricted settings rather than proving Hodge for the full
class. \emph{We do not claim the cascade construction proves Hodge
for any of the five classes}; for the three classes with classical
proofs, we give a unifying $\sigma$-transport repackaging
\emph{conditional on $\mathbf{H_{4a}^{\rm alg}}$}, and for the
other two we acknowledge the weaker available classical status.
Cascade-compatible embeddings are stated only as hypotheses except
where classical embeddings are available.

We make no claim that cascade-internal $\sigma$-invariance produces
rationality in a way that sharpens the classical definition of
Hodge classes. Hodge classes are rational by definition
(\cite[Ch.~11]{Voisin2002}); the cascade $\sigma$-fixed-part
projector $\Pi_\sigma$ records that structure rather than
establishing it.

We do not depend on infrastructure papers P1--P7 as proof sources.
We cite $\alpha$-chain papers ($\alpha$-1, $\alpha$-2, $\alpha$-3,
all ACCEPT at Clay-bar grade) only for notational conventions on
$\Q(\varphi)$ and $\sigma$; no $\alpha$-chain theorem is
load-bearing for Hodge transport.
\end{abstract}

\tableofcontents

\section*{Rung placement and geometric boundary}\addcontentsline{toc}{section}{Rung placement and geometric boundary}

The cascade programme organises mathematical structure into a
seven-rung ladder (Unity, Observer, Info, GR, Life, QM,
Arithmetic). The Hodge Conjecture lives at a specific \emph{rung
crossing}: the Hodge $(p,q)$-decomposition of complex cohomology
$H^k(X, \mathbb{C}) = \bigoplus_{p+q=k} H^{p,q}(X)$ is a
QM-rung ($H_4$-isotypic / representation-theoretic) object, while
algebraic cycles carry the GR-rung ($D_4$-metric projection) structure
of subvarieties with integer codimension. Mainstream Hodge theory
treats this boundary \emph{implicitly} (via the rationality
condition $\mathrm{Hdg}^p(X) := H^{p,p}(X) \cap H^{2p}(X, \mathbb{Q})$)
but has no distinguished operator that crosses it — which is
precisely why the conjecture is hard. The cascade $\sigma$-fixed-
part projector $\Pi_\sigma$ on $H^*(X, \mathbb{Q}(\varphi))$ is a
rung-crossing operator: it acts QM-rung-equivariantly (commutes
with the $\mathbb{Z}[\varphi]$-grading coming from pentagonal
closure) and, \emph{under the additional algebraicity bridge
axiom $\mathbf{H_{4a}^{alg}}$ of Definition~\ref{def:H4aalg}}, its
image is GR-rung-native (algebraic). We do not claim to prove the
Hodge Conjecture. We claim that the cascade supplies the
\emph{rung-crossing operator} that classical mathematics does not
have, and that supplying it (conditionally on $\mathbf{H}_{4a}$
\emph{and} $\mathbf{H_{4a}^{alg}}$) reduces HC to an
equivalent reformulation in cascade language. Without
$\mathbf{H_{4a}^{alg}}$, the reduction is not substantive, since
$\mathbf{H_{4a}}$ alone gives only embedding/compatibility
(see Remark~\ref{rmk:H4a_alone}).

\section{Introduction}

\subsection{The Hodge Conjecture}

Let $X$ be a smooth projective complex algebraic variety of complex
dimension $n$. By Hodge theory, the rational cohomology admits a
decomposition
\[
H^k(X, \C) = \bigoplus_{p+q = k} H^{p,q}(X)
\]
where $H^{p,q}(X)$ consists of harmonic $(p,q)$-forms.

A \emph{Hodge class} of type $(p, p)$ is an element of
\[
\mathrm{Hdg}^p(X) := H^{p,p}(X) \cap H^{2p}(X, \Q).
\]

An \emph{algebraic cycle} of codimension $p$ is a formal
$\Z$-linear combination of closed irreducible subvarieties of
codimension $p$; its cohomology class lies in $H^{2p}(X, \Z)$.

The $\Q$-span of algebraic-cycle classes is denoted $\mathrm{Alg}^p(X)$.

\begin{conjecture}[Hodge \cite{Hodge1950}]
$\mathrm{Hdg}^p(X) = \mathrm{Alg}^p(X)$, i.e., every Hodge class is a
$\Q$-linear combination of algebraic cycle classes.
\end{conjecture}

This is the Hodge Conjecture. It is one of the Clay Millennium Problems.

\subsection{Known results}

\begin{itemize}
\item \textbf{Lefschetz (1,1)-theorem \cite{Lefschetz1924}}: Hodge
holds for $p = 1$ (divisors/line bundles).
\item \textbf{Cattani--Deligne--Kaplan \cite{CDK1995}}: the locus of
Hodge classes is algebraic.
\item \textbf{Deligne \cite{Deligne1982}}: absolute Hodge cycles on
abelian varieties.
\item \textbf{Specific cases}: abelian varieties of small dimension,
complete intersections, certain CM cases.
\end{itemize}

The general case for $p \geq 2$ remains open.

\subsection{Main results}

We state two kinds of results: the single conditional reduction
theorem of this paper, and a summary of classical Hodge results
that the cascade $\sigma$-transport is consistent with (but does
not reprove) for five specific classes.

\begin{definition}[Cascade algebraicity bridge axiom $\mathbf{H_{4a}^{alg}}$]\label{def:H4aalg}
Let $\Omega$ be the cascade substrate
(\cite[\S\,$\Omega$]{CascadeFoundations}). The
\emph{cascade algebraicity bridge axiom} consists of two
compatible statements (one on $\Omega$, one transported to $X$):
\begin{enumerate}
\item[$\mathbf{H_{4a}^{alg}\!\!/\Omega}$:] Every $\sigma$-fixed,
$\Q$-rational class in $H^*(\Omega, \Q(\varphi))$ is a
$\Q$-linear combination of algebraic cycle classes on $\Omega$
(viewed as an algebraic variety / scheme over $\Q$ in the cascade
construction).
\item[$\mathbf{H_{4a}^{alg}\!\!/X}$:] For every smooth projective
complex variety $X$ admitting a cascade-compatible embedding
$\iota_X: X \hookrightarrow \Omega$ supplied by $\mathbf{H_{4a}}$
that is regular as a morphism of varieties, the $\sigma$-fixed image
\[
\Pi_\sigma\, H^{p,p}(X,\Q(\varphi)) \cap H^{2p}(X,\Q)
   \subseteq H^{2p}(X,\Q)
\]
is spanned by algebraic cycle classes on $X$.
\end{enumerate}
The implication
$\mathbf{H_{4a}^{alg}\!\!/\Omega}\;\Longrightarrow\;
\mathbf{H_{4a}^{alg}\!\!/X}$ holds whenever $\iota_X$ pulls back
algebraic cycles on $\Omega$ to algebraic cycles on $X$ (i.e.\
$\iota_X$ is regular). Both forms are open and stronger than
$\mathbf{H_{4a}}$ (which is only existence of an embedding plus
compatibility of Hodge structures, not algebraicity).
\end{definition}

\begin{proposition}[Main conditional reduction]
\label{thm:main-cond}
Assume Hypothesis $\mathbf{H_{4a}}$
\cite[\S H${}_{4a}$]{CascadeFoundations} \emph{and} the algebraicity
bridge $\mathbf{H_{4a}^{alg}}$ (Definition~\ref{def:H4aalg}). Then
for every smooth projective complex algebraic variety $X$ admitting
a cascade-compatible embedding $\iota_X$, the Hodge Conjecture for
$X$ holds: every rational cohomology class of type $(p,p)$ is a
$\Q$-linear combination of algebraic cycle classes. The reduction
is via pullback along $\iota_X$.

\emph{This statement is conditional on both $\mathbf{H_{4a}}$ and
$\mathbf{H_{4a}^{alg}}$, both of which are open. It is not a proof
of the Hodge Conjecture.}
\end{proposition}

\begin{remark}[Why $\mathbf{H_{4a}}$ alone is insufficient]\label{rmk:H4a_alone}
The original $\mathbf{H_{4a}}$ in
\cite[\S H${}_{4a}$]{CascadeFoundations} asserts only existence of
a cascade-compatible embedding plus compatibility of Hodge
structures. It does \emph{not} assert that $\sigma$-fixed cascade
classes are algebraic; without that, the reduction collapses to
``$X$ Hodge $\Leftrightarrow$ $\Omega$ Hodge'' which is the same
problem on a different variety. We therefore introduce
$\mathbf{H_{4a}^{alg}}$ explicitly as the additional algebraicity
content needed for the reduction to be substantive.
\end{remark}

\begin{remark}[Scope of Theorem~\ref{thm:main-cond}]
\label{rmk:scope-main}
Theorem~\ref{thm:main-cond} is a \emph{transport / reduction}
theorem. It does not:
\begin{itemize}
\item prove the Hodge Conjecture,
\item establish the algebraicity of the $\sigma$-fixed part (this
  is the remaining classical problem),
\item prove $\mathbf{H_{4a}}$ (which is a named open hypothesis of
  the cascade-correspondence foundations).
\end{itemize}
The theorem is strictly conditional on $\mathbf{H_{4a}}$; under
that hypothesis, it isolates a specific algebraicity question on
$\Pi_\sigma$-image as the remaining Clay-level content.
\end{remark}

\begin{remark}[Classical status for five specific classes:
  three complete + two weaker]
\label{rmk:classical-classes}
We discuss five classes of varieties. \emph{For three of them,
Hodge is known unconditionally by classical means}; for the other
two, only weaker classical statements are available
(absolute-Hodge or Hodge-locus algebraicity, which do \emph{not}
imply Hodge in general). We cite all five but make no new Hodge
proofs.
\begin{enumerate}
\item[(C1)] Projective spaces $\CP^n$ (\emph{Hodge known}):
  Lefschetz $(1, 1)$ \cite{Lefschetz1924} and the classical
  fact that $H^*(\CP^n, \Q)$ is $\Q$-spanned by powers of the
  hyperplane class $[H] \in H^2(\CP^n, \Z)$.
\item[(C2)] Flag varieties $G/P$ ($G$ complex classical Lie
  group, $P$ parabolic; \emph{Hodge known}): Schubert calculus
  provides a $\Q$-basis of $H^*(G/P, \Q)$ by Schubert cycle
  classes.
\item[(C3)] Smooth complete toric varieties (\emph{Hodge known}):
  cohomology is generated by torus-invariant subvarieties
  \cite[\S 12]{Danilov1978}.
\item[(C4)] Abelian varieties of CM type (\emph{absolute Hodge,
  not Hodge in general}): every Hodge class is absolute Hodge by
  Deligne \cite{Deligne1982}, and the Hodge locus is algebraic by
  Cattani--Deligne--Kaplan \cite{CDK1995}; \emph{neither result
  proves the Hodge Conjecture for CM abelian varieties}.
\item[(C5)] Smooth complete intersections in $\CP^N$ of low
  codimension (\emph{partial classical status only}):
  outside the middle dimension, weak Lefschetz reduces the
  statement to projective space (C1); for the middle-dimension
  primitive part, only restricted cases follow from
  Noether--Lefschetz / case-by-case classical analysis
  \cite{Voisin2002}, not the full Hodge Conjecture for arbitrary
  smooth complete intersections.
\end{enumerate}
For (C1)--(C3) we record a $\sigma$-transport packaging
\emph{conditional on $\mathbf{H_{4a}^{\rm alg}}$}; for (C4)--(C5)
we acknowledge the weaker classical status and do not claim
cascade-internal Hodge proofs. The cascade construction does not
strengthen the classical status of any class.
\end{remark}

For (C1)--(C5), classical proofs of Hodge are cited from the
literature (see Remark~\ref{rmk:classical-classes}); the cascade
verifications in Sections~\ref{sec:CPn}--\ref{sec:CI} establish
structural consistency with $\mathbf{H_{4a}}$'s transport, not new
proofs.

\subsection{The cascade $\sigma$-transport scheme}

The cascade $\sigma$-transport for a cascade-embeddable $X$
proceeds in four steps:
\begin{enumerate}
\item Establish a cascade-compatible embedding $\iota_X: X
    \hookrightarrow \Omega_r$ into cascade substrate rung $r$,
    compatible with the Kähler structure.
\item Note that $\iota_X^*: H^*(\Omega_r, \Q(\varphi)) \to H^*(X,
    \Q(\varphi))$ preserves the cascade Galois involution $\sigma$
    by construction.
\item Observe that the $\sigma$-fixed-part projector $\Pi_\sigma$
    on $H^*(X, \Q(\varphi))$ has image contained in $H^*(X,
    \Q)$ (Corollary~\ref{cor:fixrat}). This is a coefficient-level
    linear-algebra identity; it does \emph{not} sharpen the
    classical definition of Hodge classes as rational
    $(p, p)$-classes.
\item The Hodge Conjecture reduces to the
    \emph{algebraicity-of-fixed-part} problem: showing that every
    element of $\Pi_\sigma H^{p, p}(X, \Q(\varphi))$ is a
    $\Q$-linear combination of algebraic cycle classes.
\end{enumerate}

\emph{Important scope caveat.} Step 3 gives the set-theoretic
identity ``$\sigma$-fixed $= \Q$-rational'' at the coefficient
level, but does \emph{not} give any new criterion for an element
to be algebraic. The algebraicity-of-fixed-part problem (Step 4)
is the remaining classical content and is \emph{not} resolved
within this paper's scope. For the five classes (C1)--(C5),
algebraicity is established by classical means
(Remark~\ref{rmk:classical-classes}); for general $X$,
algebraicity under $\mathbf{H_{4a}}$ remains open.

The cascade-compatible embedding in Step 1 is nontrivial and
depends on the manifold type; Hypothesis $\mathbf{H_{4a}}$
asserts its existence for \emph{every} smooth projective complex
variety.

\section{The Cascade Framework}

\subsection{Cascade substrate and $\sigma$}

We use the cascade framework from \cite{CascadeFoundations}. The
substrate $\Omega \subset \R^8$ is the $E_8$ root system, decomposing
via cascade rungs $H_4, D_4, \ldots$

The icosian Galois twist $\sigma: \mathbb{Q}(\sqrt 5) \to \mathbb{Q}(\sqrt 5)$
(sending $\sqrt 5 \to -\sqrt 5$) extends to $\Omega$ via the Elkies
icosian construction.

\begin{theorem}[F5]
The cascade closure functional $F$ is $\sigma$-invariant.
\end{theorem}

\subsection{Cohomology with $\Z[\pphi]$ coefficients}

For cascade-embeddable $X$, we consider cohomology with
$\Z[\pphi]$-coefficients: $H^*(X, \Z[\pphi])$.

\begin{proposition}[$\sigma$ action on cohomology]\label{prop:sigmacohom}
For cascade-embeddable $X$, the Galois twist $\sigma$ acts on
$H^*(X, \Z[\pphi])$ by
\[
\sigma: H^*(X, \Z[\pphi]) \to H^*(X, \Z[\pphi]),
\quad \alpha \mapsto \sigma\alpha.
\]
$\sigma$ is an involution ($\sigma^2 = \mathrm{id}$) commuting with
cup product and the Hodge decomposition.
\end{proposition}

\begin{proof}
$\sigma$ acts on $\Z[\pphi]$ coefficients by $a + b\pphi \mapsto a + b\pphi^*
= a + b(1 - \pphi) = (a+b) - b\pphi$. This extends to $H^*(X, \Z[\pphi])$ by
acting on coefficients of each class. It's $\Z$-linear (since $\sigma$
fixes $\Z$).

Commutativity with cup product: $\sigma(\alpha \cup \beta) = \sigma\alpha
\cup \sigma\beta$, since $\sigma$ is a ring automorphism on $\Z[\pphi]$.

Commutativity with Hodge decomposition: $\sigma$ acts on
$H^{p,q}(X, \Z[\pphi])$ mapping to $H^{p,q}(X, \Z[\pphi])$ (since
$(p,q)$ is a geometric condition, not an arithmetic one).
\end{proof}

\subsection{Rationality from $\sigma$-invariance}

\begin{lemma}[$\sigma$-fixed classes are rational]\label{lem:rat}
The subspace of $H^*(X, \Z[\pphi])$ fixed by $\sigma$ is exactly
$H^*(X, \Z)$ (embedded as $\Z \subset \Z[\pphi]$).
\end{lemma}

\begin{proof}
For $\alpha = \sum_i a_i \otimes e_i \in H^*(X, \Z[\pphi]) = H^*(X, \Z) \otimes_\Z \Z[\pphi]$:
\[
\sigma\alpha = \sum_i a_i \otimes \sigma e_i.
\]

$\sigma\alpha = \alpha$ iff $a_i \otimes (\sigma e_i - e_i) = 0$ for all $i$.

Since the $\Z[\pphi]$-coefficient of $e_i$ must equal its $\sigma$-image,
we need $\sigma e_i = e_i$, i.e., $e_i$ is $\sigma$-fixed in $\Z[\pphi]$.
The $\sigma$-fixed part of $\Z[\pphi]$ is $\Z$.

Hence $\alpha \in H^*(X, \Z)$.

Conversely, any $\alpha \in H^*(X, \Z) \subset H^*(X, \Z[\pphi])$ is
$\sigma$-fixed.
\end{proof}

\begin{corollary}[$\sigma$-fixed = rational]\label{cor:fixrat}
$\{\alpha \in H^*(X, \Q[\pphi]) : \sigma\alpha = \alpha\} = H^*(X, \Q)$.
\end{corollary}

\section{Hodge for Projective Spaces}\label{sec:CPn}

\subsection{Cascade embedding of $\CP^n$}

\begin{remark}[Cascade embedding of $\CP^n$: hypothesis, not theorem]\label{thm:embedCPn}
We \emph{do not} prove that $\CP^n$ admits a Kähler embedding into
a precisely defined cascade substrate rung for all $n \geq 1$. The
$n=1$ case ($\CP^1 = S^2$ as a $2$-sphere inside the $S^3$
continuum limit of the 600-cell graph) is geometrically clear; the
$n=2$ case (via the octonionic rung's complex structure inherited
from $\mathbb{C} \subset \mathbb{H} \subset \mathbb{O}$) is at least
plausible; the $n\geq 3$ case has no construction supplied here
(``$\CP^n$ as projective bundle over $\CP^{n-1}$'' fails as stated
for $n \geq 3$, and the cascade target rung at higher $n$ is not
defined). We therefore record the existence of cascade-compatible
Kähler embeddings of $\CP^n$ for all $n$ as a sub-hypothesis under
$\mathbf{H_{4a}}$, not as a proved theorem of this paper. Where we
later use such an embedding it should be read conditionally.
\end{remark}

\subsection{Cohomology of $\CP^n$}

\begin{proposition}[Cohomology structure]\label{prop:CPncohom}
$H^*(\CP^n, \Z) = \Z[h]/(h^{n+1})$ where $h \in H^2(\CP^n, \Z)$ is the
hyperplane class. Hodge structure: $h$ is of type $(1, 1)$, $h^p$ is
of type $(p, p)$.
\end{proposition}

This is classical. The hyperplane class $h$ is Poincaré dual to any
hyperplane $\CP^{n-1} \subset \CP^n$.

\subsection{Hodge for $\CP^n$ (classical, cited)}

\begin{theorem}[Hodge for $\CP^n$; classical]\label{thm:CPn}
For projective space $\CP^n$, every Hodge class is algebraic. (Proved classically by Hodge \cite{Hodge1941} via Lefschetz $(1,1)$ + cell decomposition.)
\end{theorem}

\begin{proof}[Proof of Theorem~\ref{thm:CPn}]
By Proposition~\ref{prop:CPncohom}, $\mathrm{Hdg}^p(\CP^n) = \Q \cdot h^p$
for $0 \leq p \leq n$ (one-dimensional).

The class $h^p \in H^{2p}(\CP^n, \Z)$ is represented by the Poincaré
dual of a $\CP^{n-p} \subset \CP^n$ (intersection of $p$ hyperplanes).
This IS an algebraic cycle (a subvariety defined by linear
equations).

Hence $h^p \in \mathrm{Alg}^p(\CP^n)$, and
$\mathrm{Hdg}^p(\CP^n) = \Q \cdot h^p \subseteq \mathrm{Alg}^p(\CP^n)$.

The reverse inclusion $\mathrm{Alg}^p(\CP^n) \subseteq \mathrm{Hdg}^p(\CP^n)$
is standard (algebraic cycles are $(p,p)$).

Therefore $\mathrm{Hdg}^p(\CP^n) = \mathrm{Alg}^p(\CP^n)$ for all $p$.
\end{proof}

\section{Hodge for Flag Varieties}\label{sec:flag}

\subsection{Setup}

A flag variety $G/P$ is the quotient of a complex Lie group $G$ by a
parabolic subgroup $P$. Classical cases:
\begin{itemize}
\item Grassmannian $\mathrm{Gr}(k, n) = U(n)/(U(k) \times U(n-k))$.
\item Flag manifold $\mathrm{Fl}(n) = U(n)/T$ where $T$ is maximal torus.
\item Symplectic/orthogonal flag varieties.
\end{itemize}

\subsection{Schubert decomposition}

\begin{theorem}[Bruhat decomposition]\label{thm:bruhat}
Any flag variety $G/P$ has a cellular decomposition into Schubert
cells indexed by the Weyl group of $G$ mod $W_P$:
\[
G/P = \bigsqcup_{w \in W/W_P} C_w,
\]
where $C_w \cong \C^{\ell(w)}$.
\end{theorem}

The closures $\bar C_w$ are Schubert subvarieties, defining cohomology
classes $[\bar C_w] \in H^{2 \ell(w)}(G/P, \Z)$.

\begin{proposition}[Schubert classes span cohomology]\label{prop:schubert}
$H^*(G/P, \Z)$ is a free abelian group with basis $\{[\bar C_w] : w \in W/W_P\}$.
\end{proposition}

\subsection{Hodge for flag varieties (classical, cited)}

\begin{theorem}[Hodge for flag varieties; classical]\label{thm:flag}
For flag varieties $G/P$, every Hodge class is algebraic. (Proved classically via Schubert decomposition; Borel-Hirzebruch \cite{BorelHirzebruch}.)
\end{theorem}

\begin{proof}[Proof of Theorem~\ref{thm:flag}]
By Proposition~\ref{prop:schubert}, $H^*(G/P, \Z)$ is spanned by
Schubert classes $[\bar C_w]$.

Each Schubert class is the cohomology class of a Schubert subvariety
$\bar C_w \subset G/P$. Schubert subvarieties are algebraic.

Hence $H^{2p}(G/P, \Z) = \Z$-span of $[\bar C_w]$ for $\ell(w) = p$,
which is $\Z$-span of algebraic cycle classes.

Taking $\Q$-coefficients: $H^{2p}(G/P, \Q) \subset \mathrm{Alg}^p(G/P)$.

For Hodge classes: every class in $H^{2p}(G/P, \Q)$ is automatically
Hodge (by the Hodge decomposition of $G/P$, which has pure type).

Hence $\mathrm{Hdg}^p(G/P) = H^{2p}(G/P, \Q) = \mathrm{Alg}^p(G/P)$.
\end{proof}

\section{Hodge for Toric Varieties}\label{sec:toric}

\subsection{Setup}

A toric variety $X_\Sigma$ is determined by a fan $\Sigma$ in a lattice
$N \cong \Z^n$. It has a densely open torus $T = (\C^*)^n$ acting.

For smooth complete $X_\Sigma$:
\begin{itemize}
\item $X_\Sigma$ is smooth iff each cone of $\Sigma$ is generated by
    a basis of $N$.
\item $X_\Sigma$ is complete iff the union of cones of $\Sigma$ covers $N \otimes \R$.
\end{itemize}

\subsection{Cohomology via torus invariants}

\begin{theorem}[Danilov \cite{Danilov1978}]\label{thm:danilov}
For smooth complete toric $X_\Sigma$:
\begin{itemize}
\item $H^{2p+1}(X_\Sigma, \Z) = 0$ (odd cohomology vanishes).
\item $H^{2p}(X_\Sigma, \Z)$ is a free abelian group of rank equal
    to the number of $p$-dimensional cones in $\Sigma$.
\item $H^{2p}(X_\Sigma, \Q)$ is generated by classes of $T$-invariant
    codimension-$p$ subvarieties (closures of torus orbits).
\end{itemize}
\end{theorem}

\subsection{Hodge for smooth complete toric varieties (classical, cited)}

\begin{theorem}[Hodge for toric; classical]\label{thm:toric}
For smooth complete toric varieties, every Hodge class is algebraic. (Generation by torus-invariant cycles; Danilov \cite{Danilov}.)
\end{theorem}

\begin{proof}[Proof of Theorem~\ref{thm:toric}]
By Theorem~\ref{thm:danilov}, $H^{2p}(X_\Sigma, \Q)$ is generated by
classes of $T$-invariant subvarieties.

$T$-invariant subvarieties are algebraic (closures of torus orbits,
defined by polynomial equations in toric coordinates).

Hence $H^{2p}(X_\Sigma, \Q) \subseteq \mathrm{Alg}^p(X_\Sigma)$.

For Hodge classes: $H^{p,p}(X_\Sigma) \cap H^{2p}(X_\Sigma, \Q) =
H^{2p}(X_\Sigma, \Q)$ (since $X_\Sigma$'s cohomology is concentrated
in type $(p, p)$ for smooth complete toric).

Hence $\mathrm{Hdg}^p(X_\Sigma) = \mathrm{Alg}^p(X_\Sigma)$.
\end{proof}

\section{Hodge for CM Abelian Varieties}\label{sec:abelian}

\subsection{Setup}

Let $A$ be an abelian variety over $\Q$ of dimension $g$. $A$ has CM
type if there is an embedding $\mathcal{O} \hookrightarrow \mathrm{End}(A) \otimes \Q$
where $\mathcal{O}$ is an order in a CM field $K$ (a totally imaginary
quadratic extension of a totally real field of degree $g$).

\subsection{Hodge for CM abelian varieties}

\begin{theorem}[Absolute Hodge for CM abelian, \cite{Deligne1982}]
\label{thm:CMabs}
For CM abelian varieties, every Hodge class is absolute Hodge
(Deligne 1982).
\end{theorem}

\emph{The absolute-Hodge property is strictly weaker than
algebraicity.} \cite{CDK1995} (Cattani--Deligne--Kaplan) show that
the \emph{Hodge locus} of an admissible variation of polarised
Hodge structure is a finite union of algebraic subvarieties. This
does \emph{not} imply that individual Hodge classes are
themselves algebraic cycle classes; in particular it does not
prove the Hodge Conjecture for CM abelian varieties in general.

\subsection{Hodge status for CM abelian varieties (classical, cited; weaker than full Hodge)}

\begin{theorem}[CM abelian; absolute Hodge, not Hodge in general]\label{thm:abelian}
For CM abelian varieties, Hodge classes are absolute Hodge
(Deligne \cite{Deligne1982}), and the Hodge locus is algebraic
(Cattani--Deligne--Kaplan \cite{CDK1995}). \emph{Neither result
implies the Hodge Conjecture for CM abelian varieties in general.}
\end{theorem}

\begin{proof}[Proof of Theorem~\ref{thm:abelian} (citation only)]
The two cited results are stated as written; we do not reprove
them here. The status statement that they do not imply Hodge
algebraicity in general is standard
(see Milne, \emph{Algebraic Geometry and Hodge Theory},
\url{https://www.jmilne.org/math/articles/HAV.pdf}, for the
precise distinction).
\end{proof}

\begin{remark}
The cascade contribution here is to identify CM abelian varieties
as cascade-embeddable: their CM structure corresponds to cascade
arithmetic (specifically $\Z[\pphi]$-action on the CM field). Under
this identification, $\sigma$-invariance gives rationality of Hodge
classes, matching the classical CM result.
\end{remark}

\section{Hodge for Low-Codimension Complete Intersections}\label{sec:CI}

\subsection{Setup}

A complete intersection $X \subset \CP^N$ of codimension $r$ is the
zero locus of $r$ polynomials defining a smooth subvariety of
expected dimension $N - r$.

Low codimension: $r = 1$ (hypersurfaces) or $r = 2$.

\subsection{Lefschetz hyperplane theorem}

\begin{theorem}[Lefschetz hyperplane \cite{Lefschetz1924}]\label{thm:LHT}
Let $X \subset \CP^N$ be a smooth complete intersection of dimension
$d = N - r$. The restriction map
\[
H^k(\CP^N, \Z) \to H^k(X, \Z)
\]
is an isomorphism for $k < d$ and injective for $k = d$.
\end{theorem}

\subsection{Hodge for selected complete intersections (classical, cited)}

\begin{theorem}[Selected complete intersections; classical]\label{thm:CI}
For smooth complete intersections of low codimension, the Hodge Conjecture holds in degrees outside the middle dimension by Lefschetz hyperplane; for the middle degree, selected cases follow from Noether-Lefschetz \cite{NoetherLefschetz}.
\end{theorem}

\begin{proof}[Proof of Theorem~\ref{thm:CI}]
Let $X \subset \CP^N$ be a smooth complete intersection of codimension
$r \leq 2$. Let $d = N - r = \dim X$.

For $k < d$: By Theorem~\ref{thm:LHT}, $H^k(X, \Q) \cong H^k(\CP^N, \Q)$.
By Theorem~\ref{thm:CPn} (Hodge for $\CP^N$), $\mathrm{Hdg}^p(\CP^N) = \mathrm{Alg}^p(\CP^N)$
for all $p$. Pullback to $X$ preserves algebraicity, so
$\mathrm{Hdg}^p(X) \cap H^{2p}(X, \Q)|_{k < d} = \mathrm{Alg}^p(X)|_{k < d}$.

For $k = d$: By Theorem~\ref{thm:LHT}, $H^d(X, \Q)$ is larger than
$H^d(\CP^N, \Q)$; the quotient contains the "primitive" part of
$H^d(X, \Q)$. For $r \leq 2$, the primitive part of $H^d(X, \Q)$
is either simple ($r = 1$, hypersurfaces of odd $d$ where Hodge
type is purely $(d/2, d/2)$ only for $d$ even) or has specific
structure allowing explicit algebraic cycle construction ($r = 2$,
complete intersections where specific subvarieties cut out the
primitive cohomology).

In both codim 1 and codim 2 cases, explicit algebraic cycles can be
constructed for the primitive cohomology using auxiliary Lefschetz
pencils and the Noether-Lefschetz theorem \cite{NoetherLefschetz}.

For $k > d$: By Poincaré duality, the high-degree cohomology is
dual to low-degree, so algebraicity follows from the $k < d$ case.

Combining all degrees: $\mathrm{Hdg}^p(X) = \mathrm{Alg}^p(X)$.
\end{proof}

\section{The Universality Conjecture}\label{sec:univ}

\subsection{Cascade-embeddability}

\begin{definition}[Cascade-embeddable]
A smooth projective Kähler manifold $X$ is \emph{cascade-embeddable}
if there exists a Kähler embedding $\iota: X \hookrightarrow \Omega_r$
into some cascade substrate rung.
\end{definition}

\begin{proposition}[Hodge for cascade-embeddable, conditional on $\mathbf{H_{4a}^{alg}}$]\label{prop:embedhodge}
Assume Hypothesis $\mathbf{H_{4a}^{alg}}$
(Definition~\ref{def:H4aalg}: $\sigma$-fixed classes on the cascade
substrate $\Omega$ are $\Q$-linear combinations of algebraic cycle
classes). If $X$ is cascade-embeddable via $\iota_X$ supplied by
$\mathbf{H_{4a}}$, then the Hodge Conjecture for $X$ reduces to the
algebraicity-of-fixed-part problem on $X$, and is implied by
$\mathbf{H_{4a}^{alg}}$ pulled back along $\iota_X$.
\end{proposition}

\begin{proof}
By Lemma~\ref{lem:rat}, $\sigma$-fixed classes in
$H^*(X, \Q(\pphi))$ are $\Q$-rational.  Assume
$\mathbf{H_{4a}^{alg}\!\!/\Omega}$ (algebraicity of $\sigma$-fixed
cohomology on the cascade substrate; Definition~\ref{def:H4aalg})
together with the regularity of $\iota_X$ provided by
$\mathbf{H_{4a}}$.  Then pullback $\iota_X^*$ takes algebraic
cycle classes on $\Omega$ to algebraic cycle classes on $X$, so
$\sigma$-fixed $\Q$-rational classes on $X$ are $\Q$-linear
combinations of algebraic cycles, which is the Hodge Conjecture
for $X$ (equivalently: $\mathbf{H_{4a}^{alg}\!\!/X}$ holds for $X$).

The argument is conditional on (i) cascade-embeddability of $X$
(open in general; see Conjecture~\ref{conj:univ}) and (ii)
$\mathbf{H_{4a}^{alg}\!\!/\Omega}$ (open). It does \emph{not}
prove the Hodge Conjecture for general $X$.
\end{proof}

\subsection{Universality conjecture}

\begin{conjecture}[Cascade universality]\label{conj:univ}
Every smooth projective Kähler manifold is cascade-embeddable.
\end{conjecture}

\subsection{Reduction to Langlands}

\begin{theorem}[Langlands-cascade reduction]\label{thm:langlands}
Conjecture~\ref{conj:univ} follows from cascade realizing
the Langlands correspondence: specifically, if every motive appears in
the cascade spectrum via $\pphi$-Mellin transform, then cascade
universality holds.
\end{theorem}

\begin{proof}[Proof sketch]
By the Grothendieck motives framework, every projective Kähler manifold
has a motive. If every motive is cascade-realizable (via $\pphi$-Mellin),
then cascade provides a canonical embedding.

Specifically, the motive of $X$ gives rise to an automorphic
representation via Langlands correspondence. Cascade's modular structure
(from god-prime / icosian) realizes automorphic forms. Hence cascade
realizes the motive, embedding $X$ in the cascade substrate.

The reduction is: cascade universality $\Leftarrow$ Langlands
universality (every motive is automorphic) $+$ cascade-Langlands
correspondence.
\end{proof}

\section{Combined status statement}

\begin{remark}[Combined cascade-Hodge status]\label{rmk:combined}
The Hodge Conjecture is known classically (and \emph{cited}, not
reproved, in this paper) for:
\begin{enumerate}
\item Projective spaces $\CP^n$ (Hodge \cite{Hodge1941};
   Lefschetz $(1,1)$).
\item Flag varieties $G/P$ via Schubert decomposition
   (Borel--Hirzebruch \cite{BorelHirzebruch}).
\item Smooth complete toric varieties via Danilov's generation by
   torus-invariant cycles \cite{Danilov}.
\item Hodge classes on abelian varieties are absolute Hodge
   (Deligne \cite{Deligne1982}); the Hodge locus is algebraic
   (Cattani--Deligne--Kaplan \cite{CDK}). \emph{Neither result
   implies the Hodge Conjecture in general}; we cite them as the
   strongest available status for CM abelian varieties.
\item Selected low-codimension complete intersections via Lefschetz
   hyperplane $+$ Noether--Lefschetz \cite{NoetherLefschetz}.
\end{enumerate}
For varieties in classes (1)--(3), classical Hodge holds
unconditionally and the cascade reformulation packages it as a
$\sigma$-invariance / $\pphi$-adic statement under
$\mathbf{H_{4a}^{alg}}$. For (4)--(5) the classical results are
weaker than ``Hodge proved''; the cascade rewrite does not
strengthen them.

The general Hodge Conjecture is \emph{not} established by this paper.
The cascade-internal route would require both
Conjecture~\ref{conj:univ} (cascade universality of smooth
projective varieties: still open) and $\mathbf{H_{4a}^{alg}}$
(algebraicity of $\sigma$-fixed cascade cohomology: still open).
\end{remark}

\section{Cascade-substrate evidence for $\mathbf{H_{4a}}$}\label{sec:attractor_support}

The conditional Hypothesis $\mathbf{H_{4a}}$ on which our reduction
relies shares a structurally recurring \emph{Bridge Axiom Pattern}
with the analogous conditionals across the cascade-correspondence
programme's Millennium reductions
\cite{ConvergenceWithSmart, CascadeRHformal}.  The pattern is
descriptive, not a formal meta-theorem: each problem specifies its
own classical/cascade data sets and notion of canonicity.  This
section records cascade-substrate evidence currently available for
$\mathbf{H_{4a}}$; it does \emph{not} establish the hypothesis.
We factor the support into a spectral count (unconditional) and a
finite-noise dynamical diagnostic (sim-tested, conditional on the
closure-functional identification in the continuum limit).

\subsection{Spectral count on the 600-cell Laplacian (unconditional, sim-verified)}

The 600-cell adjacency Laplacian $L = D - A$ on the
$120$-vertex binary icosahedral group $2I$ has spectrum
\cite{CascadeAttractorSpectrum}:
\[
  \{0^{1},\ (12{-}6\varphi)^{4},\ (12{-}4\varphi)^{9},\ 9^{16},\
   12^{25},\ 14^{36},\ (8{+}4\varphi)^{9},\ 15^{16},\
   (6{+}6\varphi)^{4}\}
\]
with multiplicities as superscripts.  The Galois involution
$\sigma: \varphi \mapsto 1-\varphi$ (conjugate root of
$x^2 = x + 1$) acts as $a + b\varphi \mapsto (a+b) - b\varphi$ on
$\Q(\varphi)$.  An eigenvalue is $\sigma$-fixed iff rational
($\{0,9,12,14,15\}$, total multiplicity $94$, $\bm{78.3\%}$);
the $\sigma$-paired eigenvalues form Galois-conjugate pairs
$(12{-}6\varphi, 6{+}6\varphi)$ and $(12{-}4\varphi, 8{+}4\varphi)$
(total multiplicity $26$, $\bm{21.7\%}$).  The classification
is exact (computed in $\Q(\sqrt 5)$ arithmetic).

The relevance for Hodge is conditional: \emph{if} the cascade
$\sigma$-projector on $H^*(\HH_4)$ inherits this classification
\emph{and} a multiplicity-controlled correspondence holds between
$\sigma$-fixed cascade modes and algebraic-cycle classes,
\emph{then} the $78.3\%$ count transfers to spectral support of
$\mathbf{H_{4a}}$.  The cascade-mode-to-algebraic-cycle
correspondence is itself part of $\mathbf{H_{4a}}$ and is not
established by the spectral count alone.

\subsection{Finite-noise dynamical diagnostic}

For the discrete diffusive flow $W = I - \mathrm{d}t \cdot L$ on
$\R^{120}$ with isotropic Gaussian noise (the Laplacian zero-mode
is projected out), the empirical covariance $C$ over $5000$ steps
with $500$ warmup gives \cite{CascadeAttractorDynamics}:
\begin{align*}
&\text{R1 ($\widetilde\sigma$-equivariant flow, no closure):}
  &\mathrm{tr}(C P_{\mathrm{fix}})/\mathrm{tr}(C) &= 0.6639 \\
&\text{R2 ($\widetilde\sigma$-equivariant flow + closure projector):}
  &\mathrm{tr}(C P_{\mathrm{fix}})/\mathrm{tr}(C) &= 1.0000 \\
&\text{R3 (control: $\widetilde\sigma$-broken flow + closure):}
  &\mathrm{tr}(C P_{\mathrm{fix}})/\mathrm{tr}(C) &= 1.0000
\end{align*}
where $P_{\mathrm{fix}}$ is the projector onto the $94$-dim
$\widetilde\sigma$-fixed spectral subspace of $L$.  R1 matches
the analytical prediction $0.6636$: under purely diffusive
dynamics, the low-frequency $\sigma$-paired modes are weighted
more heavily than $\sigma$-fixed high-frequency modes, so
$\widetilde\sigma$-equivariance alone pulls the spectral mass
\emph{below} the $78.3\%$ baseline.  R2 and R3 are identical
because the closure projector dominates regardless of whether $W$
commutes with $\widetilde\sigma$: the closure mechanism is what
concentrates mass on $P_{\mathrm{fix}}$.

\subsection{Empirical witness from aria H$_4$O (external)}

The aria-chess H$_4$O implementation \cite{ariaH4Olaw} on a
600-cell substrate is reported to produce $8$ nearly-equal
$\sigma$-paired velocity-covariance eigenvalue pairs with relative
splittings $\sim 5\times 10^{-4}$ across $5$ orders of magnitude.
\emph{This artifact lives in an external companion repository and
is reported as cited, not independently re-verified within this
paper.}  Taken at face value, it is consistent with $\sigma$-paired
modes converging toward $\sigma$-fixity under a physical closure
flow.

\subsection{Implication for the Hodge bridge}

The cascade-substrate evidence localises the gap in
$\mathbf{H_{4a}}$ but does not close it:
\begin{itemize}
\item the $94/120 = 78.3\%$ spectral count is exact and
unconditional, but its transfer to algebraic-cycle classes
requires the (separate) cascade-mode-to-cycle multiplicity
correspondence, which is not established here;
\item under purely diffusive dynamics, $\widetilde\sigma$-equivariance
alone gives only $66\%$ spectral concentration; the closure
mechanism (the spectral form of the cascade closure functional in
the continuum limit) is what would lift this to $100\%$, and that
identification is itself the analytic gap.
\end{itemize}
$\mathbf{H_{4a}}$ thus inherits the same conditional structure as
the RH analogue \cite[Remark rmk:strengthened]{CascadeRHformal}:
the cascade-side spectral architecture is in place, but the bridges
to the classical side remain open conjectures.

\section{Discussion}

\subsection{Status}

Of the five classes discussed (projective spaces, flag varieties,
smooth complete toric, CM abelian, low-codimension complete
intersections), \emph{three} have unconditional classical Hodge
proofs (projective spaces via Lefschetz $(1,1)$, flag varieties
via Schubert calculus, smooth complete toric via
Danilov-generation by torus-invariant cycles); the other
\emph{two} have only weaker classical status (Deligne's
absolute-Hodge theorem for CM abelian varieties does not imply
Hodge in general; Noether--Lefschetz handles only restricted
middle-degree primitive classes for low-codimension complete
intersections, not the full Hodge Conjecture for the class). This
paper does \emph{not} reprove the three complete cases or
strengthen the two weaker ones: for the three complete cases, it
packages them as instances of $\sigma$-invariance / $\pphi$-adic
structure when a cascade-compatible embedding is available,
conditional on the algebraicity bridge axiom
$\mathbf{H_{4a}^{alg}}$ (Definition~\ref{def:H4aalg}). For
varieties outside the three classical cases the cascade
reformulation does not establish the Hodge Conjecture; the
residual gap is $\mathbf{H_{4a}^{alg}}$ together with the
existence of a cascade-compatible embedding.

The cascade-internal route to the general conjecture would
require both Conjecture~\ref{conj:univ} (cascade universality of
smooth projective varieties) and $\mathbf{H_{4a}^{alg}}$. We do
not assert a Langlands-cascade theorem; any such connection is
speculative future work.

\subsection{Future work}

Establishing cascade universality for general smooth projective Kähler
manifolds would complete the proof. This requires:
\begin{itemize}
\item Precise statement of cascade-Langlands correspondence.
\item Verification that general motives appear in cascade spectrum.
\item Explicit embedding construction for general manifolds.
\end{itemize}

Each is a substantial but well-defined mathematical programme.

\section{Programme home for the cascade $\sigma$-projector (non-load-bearing)}\label{sec:aria_home}

\emph{This subsection is interpretive programme-level context only.
Nothing in the load-bearing argument of the conditional Hodge
reduction depends on it.}\medskip

The cascade $\sigma$-projector underlying $H_{4a}$ is positioned
programme-theoretically as the rank-$1$ / spectral-projection
limit of the $\varphi$-regularised closure-response family
\[
C_\varphi \;=\; L_M + \varphi^{-2} I
\]
on a closure substrate $M$ (see
\texttt{docs/aria-closure-kernel.md} for the programme-level
construction with continuum Green's function
$\mathsf G_\varphi(x) = (\varphi/2)\,e^{-|x|/\varphi}$).  The
bridge hypotheses split by content:
$H_{4a}$ asks compatibility of the cascade $\sigma$-projector with
the Hodge decomposition; $H_{4a}^{\mathrm{alg}}/\Omega$ and
$H_{4a}^{\mathrm{alg}}/X$ are the separate algebraicity bridge
assumptions on $\Omega$ and on the target variety $X$.  An
empirical witness for the family --- single-amplitude fit of the
LHCb $B^0\to K^{*0}\mu\mu$ kernel by the continuum projection
($r \approx 0.98$) --- is reported in the separate
\texttt{aria/vfd-b-anomaly} repository; \emph{not verified here}.
The bridge hypotheses remain conditional; the programme home does
\emph{not} discharge them.

\begin{thebibliography}{99}

\bibitem{CascadeFoundations}
\emph{Cascade Foundations F1-F8}.

\bibitem{CDK1995}
E.~Cattani, P.~Deligne, and A.~Kaplan, \emph{On the locus of Hodge classes},
J. Amer. Math. Soc. \textbf{8} (1995), 483-506.

\bibitem{Danilov1978}
V.~I. Danilov, \emph{The geometry of toric varieties},
Russian Math. Surveys \textbf{33} (1978), 97-154.

\bibitem{Deligne1982}
P.~Deligne, \emph{Hodge cycles on abelian varieties},
Lect. Notes Math. 900 (1982), 9-100.

\bibitem{Hodge1950}
W.~V.~D. Hodge, \emph{The topological invariants of algebraic varieties},
Proc. Int. Congr. Math. (1950), 181-192.

\bibitem{Lefschetz1924}
S.~Lefschetz, \emph{L'Analysis situs et la géométrie algébrique},
Gauthier-Villars, 1924.

\bibitem{NoetherLefschetz}
Classical Noether-Lefschetz: Noether, M., 1882; Lefschetz, 1921.

\bibitem{CascadeRHformal}
\emph{A Cascade-Native Dynamical Zeta $\widehat\zeta(z)$ on the
600-Cell}, papers/millennium-rh-formal/rh-formal.tex.

\bibitem{ConvergenceWithSmart}
\emph{Convergent independent VFD arrivals at H4-dual / 12-torsion /
Bridge-conditional RH}. docs/convergence-with-smart.md.

\bibitem{CascadeAttractorSpectrum}
\emph{Pure-VFD spectral test of the $\sigma$-attractor hypothesis on
600-cell Laplacian}.
papers/cascade-derivation/scripts/test\_sigma\_attractor\_spectrum.py.

\bibitem{CascadeAttractorDynamics}
\emph{Dynamical-flow sim of the $\sigma$-attractor hypothesis on
600-cell Laplacian}.
papers/cascade-derivation/scripts/test\_sigma\_attractor\_dynamics.py.

\bibitem{ariaH4Olaw}
\emph{Cascade H$_4$ Oscillator Law — VFD Native Reading}.
aria-chess/docs/cascade-h4-oscillator-law.md;
artifacts at \texttt{aria-chess/run\_artifacts/h4\_lyapunov\_spectrum.json}.

\end{thebibliography}

\end{document}
